\documentclass[UTF8,11pt,a4paper]{ctexart}

\usepackage[top=1.35cm,bottom=1.35cm,left=1.7cm,right=1.7cm,footskip=0.78cm]{geometry}
\usepackage{amsmath,amssymb}
\usepackage{graphicx}
\usepackage{wrapfig}
\usepackage{xcolor}
\usepackage{calc}
\usepackage{longtable,array,tabularx}
\usepackage{enumitem}
\usepackage{eso-pic}
\usepackage{lastpage,refcount}
\usepackage{hyperref}

\definecolor{ExamInk}{HTML}{111111}
\definecolor{ExamMuted}{HTML}{4B5563}
\definecolor{ExamWatermark}{HTML}{FAFAFA}

\hypersetup{hidelinks,pdfborder={0 0 0}}
\setlength{\parindent}{0pt}
\setlength{\parskip}{0.18em}
\setlength{\emergencystretch}{3em}
\renewcommand{\arraystretch}{1.12}
\linespread{1.14}
\setcounter{secnumdepth}{0}

\setlist[enumerate,1]{
  label=\arabic*.,
  leftmargin=2.15em,
  labelsep=0.55em,
  itemsep=0.46em,
  topsep=0.32em,
  parsep=0.08em,
  partopsep=0pt
}
\setlist[enumerate,2]{leftmargin=2.1em,itemsep=0.2em,topsep=0.2em}

\providecommand{\toprule}{\hline}
\providecommand{\midrule}{\hline}
\providecommand{\bottomrule}{\hline}
\providecommand{\real}[1]{#1}

\ctexset{
  section={format=\centering\heiti\bfseries\fontsize{16pt}{22pt}\selectfont,beforeskip=0.7em,afterskip=0.55em},
  subsection={format=\heiti\bfseries\fontsize{12.5pt}{18pt}\selectfont,beforeskip=0.7em,afterskip=0.35em},
  subsubsection={format=\heiti\bfseries\normalsize,beforeskip=0.55em,afterskip=0.25em}
}

\newif\ifExamAnswers
\ExamAnswersfalse
\newcounter{ExamPaperPageCount}

\newcommand{\ExamMainTitle}[1]{%
  \begin{center}
    {\heiti\bfseries\fontsize{16pt}{23pt}\selectfont #1\par}
  \end{center}
  \vspace{0.15em}%
}
\newcommand{\ExamSubjectTitle}[1]{%
  \begin{center}
    {\heiti\bfseries\fontsize{22pt}{28pt}\selectfont #1\par}
  \end{center}
  \vspace{0.45em}%
}
\newcommand{\ExamAnswerMainTitle}[1]{%
  \vspace*{0.65cm}%
  \begin{center}
    {\heiti\bfseries\fontsize{15pt}{22pt}\selectfont #1\par}
  \end{center}
  \vspace{0.1em}%
}
\newcommand{\ExamAnswerSubjectTitle}[1]{%
  \begin{center}
    {\heiti\bfseries\fontsize{19pt}{26pt}\selectfont #1\par}
  \end{center}
  \vspace{0.7em}%
}

\newcommand{\beginExamAnswers}{%
  \setcounter{ExamPaperPageCount}{\value{page}}%
  \addtocounter{ExamPaperPageCount}{-2}%
  \refstepcounter{ExamPaperPageCount}%
  \label{ExamPaperPageTotal}%
  \clearpage
  \global\ExamAnswerstrue
  \setcounter{page}{1}%
  \fontsize{10.5pt}{21pt}\selectfont
  \setlength{\parskip}{0.24em}%
  \setlength{\emergencystretch}{6em}%
  \AddToShipoutPictureBG*{\ExamPageOverlay}%
}
\newcommand{\ExamQuestionPageBreak}{\clearpage\useExamPageStyle}
\newcommand{\ExamAnswerPageBreak}{\clearpage\useExamPageStyle}

\newcommand{\ExamFooterText}{%
  \ifExamAnswers
    第\ \thepage\ 页\quad（共\ \getpagerefnumber{LastPage}\ 页）%
  \else
    数学试卷\quad 第\ \thepage\ 页（共\ \getrefnumber{ExamPaperPageTotal}\ 页）%
  \fi
}
\newcommand{\ExamPageOverlay}{%
  \AtPageCenter{%
    \makebox(0,0){\rotatebox{42}{\textcolor{ExamWatermark}{\heiti\bfseries\fontsize{48pt}{54pt}\selectfont 数学试卷}}}}%
  \AtPageLowerLeft{%
    \raisebox{0.62cm}[0pt][0pt]{\hspace*{1.7cm}%
      \makebox[\dimexpr\paperwidth-3.4cm\relax][c]{\small\color{ExamMuted}\ExamFooterText}}}%
}
\AddToShipoutPictureBG{\ExamPageOverlay}
\pagestyle{empty}
\makeatletter
\let\ps@plain\ps@empty
\makeatother
\newcommand{\useExamPageStyle}{\AddToShipoutPictureBG*{\ExamPageOverlay}}

\newcommand{\ExamChoicesFour}[4]{%
  \par\smallskip\noindent
  \begin{tabularx}{\linewidth}{@{}>{\raggedright\arraybackslash}X>{\raggedright\arraybackslash}X>{\raggedright\arraybackslash}X>{\raggedright\arraybackslash}X@{}}
    #1 & #2 & #3 & #4
  \end{tabularx}\par\smallskip
}
\newcommand{\ExamChoicesTwo}[4]{%
  \par\smallskip\noindent
  \begin{tabularx}{\linewidth}{@{}>{\raggedright\arraybackslash}X>{\raggedright\arraybackslash}X@{}}
    #1 & #2 \\
    #3 & #4
  \end{tabularx}\par\smallskip
}
\newcommand{\ExamChoicesOne}[4]{%
  \par\smallskip\noindent #1\par
  \noindent #2\par
  \noindent #3\par
  \noindent #4\par\smallskip
}
\newcommand{\ExamChoicesFigure}[4]{%
  \par\smallskip\noindent
  \begin{minipage}[t]{0.70\linewidth}
    \noindent #1\par
    \noindent #2\par
    \noindent #3\par
    \noindent #4
  \end{minipage}\par\smallskip
}
\newcommand{\ExamScore}[1]{%
  \unskip\allowbreak\hfill\makebox[0.26\linewidth][r]{\dotfill\ #1\ 分}\par
}

\newenvironment{examSolutionBox}{%
  \par\smallskip\begingroup
  \setlength{\leftskip}{0.5em}\setlength{\rightskip}{0.5em}%
  \noindent\heiti\bfseries 解析\par\smallskip\normalfont
}{%
  \par\endgroup\smallskip
}

\makeatletter
\def\maxwidth{\linewidth}
\def\maxheight{0.12\textheight}
\setkeys{Gin}{width=\maxwidth,height=\maxheight,keepaspectratio}
\providecommand{\pandocbounded}[1]{%
  \begin{wrapfigure}{r}{0.25\textwidth}
    \vspace{-1.15em}\centering #1\vspace{-1.0em}
  \end{wrapfigure}%
}
\makeatother

\providecommand{\tightlist}{%
  \setlength{\itemsep}{0pt}\setlength{\parskip}{0pt}}


\begin{document}
\color{ExamInk}


绝密★启用前

试卷类型:A

2025---2026 学年度第二学期高二年级期末教学质量监测试卷

\ExamMainTitle{数学}

\subsection{注意事项}\label{ux6ce8ux610fux4e8bux9879}

\begin{enumerate}
\def\labelenumi{\arabic{enumi}.}
\item
  考生答卷前,务必将自己的姓名、座位号写在答题卡上。将条形码粘贴在规定区域。本试卷满分 150 分,考试时间 120 分钟。
\item
  做选择题时, 选出每小题答案后, 用铅笔把答题卡上对应题目的答案标号涂黑。如需改动, 用橡皮擦干净后, 再选涂其他答案标号。写在本试卷上无效。
\item
  回答非选择题时, 将答案写在答题卡的规定区域内, 写在本试卷上无效。
\item
  考试结束后, 将答题卡交回。
\end{enumerate}

\subsection{一、选择题:本题共8小题,每小题5分,共40分.在每小题给出的四个选项中,只有一项是符合题目要求的.}\label{ux4e00ux9009ux62e9ux9898ux672cux9898ux51718ux5c0fux9898ux6bcfux5c0fux98985ux5206ux517140ux5206.ux5728ux6bcfux5c0fux9898ux7ed9ux51faux7684ux56dbux4e2aux9009ux9879ux4e2dux53eaux6709ux4e00ux9879ux662fux7b26ux5408ux9898ux76eeux8981ux6c42ux7684.}

\begin{enumerate}
\def\labelenumi{\arabic{enumi}.}
\item
  在复平面内, 复数 \(z\) 对应的点的坐标是 \((-1, -\sqrt{3})\) , 则 \(z\) 的共轭复数 \(z =\)

  \begin{itemize}
  \tightlist
  \item
    A. \(1 + \sqrt{3}\mathrm{i}\)
  \item
    B. \(1 - \sqrt{3}\mathrm{i}\)
  \item
    C. \(-1 + \sqrt{3}\mathrm{i}\)
  \item
    D. \(-1 - \sqrt{3}\mathrm{i}\)
  \end{itemize}
\item
  已知集合 \(A = \{-1, 0, 1, 2, 8\}\) , \(B = \{x \mid \sqrt[3]{x} = x\}\) , 则 \(A \cap B =\)

  \begin{itemize}
  \tightlist
  \item
    A. \(\{0, 1\}\)
  \item
    B. \(\{-1, 0, 1\}\)
  \item
    C. \(\{0, 2, 8\}\)
  \item
    D. \(\{0, 1, 2, 8\}\)
  \end{itemize}
\item
  设 \(\vec{a}, \vec{b}\) 为单位向量，且 \(|\vec{a} - \vec{b}| = \sqrt{2}\) ，则 \(|\vec{a} + \vec{b}| =\)

  \begin{itemize}
  \tightlist
  \item
    A. 1
  \item
    B. \(\sqrt{2}\)
  \item
    C. \(\sqrt{3}\)
  \item
    D. 2
  \end{itemize}
\item
  已知 \(\dfrac{\sin(\alpha+\beta)}{\sin(\alpha-\beta)}=3\) ，则 \(\dfrac{\tan\alpha}{\tan\beta}=\)

  \begin{itemize}
  \tightlist
  \item
    A. \(\dfrac{1}{3}\)
  \item
    B. \(\dfrac{1}{2}\)
  \item
    C. 2
  \item
    D. 3
  \end{itemize}
\end{enumerate}

\ExamQuestionPageBreak

\begin{enumerate}
\def\labelenumi{\arabic{enumi}.}
\setcounter{enumi}{4}
\item
  某中学举办田径运动会, 某班从甲、乙等6名学生中选4名学生代表班级参加学校 \(4 \times 1\) 米接力赛, 其中甲只能跑第1棒或第2棒, 乙只能跑第2棒或第4棒, 那么甲、乙都参加不同棒次安排方案总数为

  \begin{itemize}
  \tightlist
  \item
    A. 12
  \item
    B. 24
  \item
    C. 30
  \item
    D. 36
  \end{itemize}
\item
  长方体的一条体对角线与它在一个顶点处的三个面所成的角分别为 \(\alpha, \beta, \gamma\) ，则

  \begin{itemize}
  \tightlist
  \item
    A. \(\sin^2\alpha + \sin^2\beta + \sin^2\gamma = 1\)
  \item
    B. \(\sin^2\alpha + \sin^2\beta + \sin^2\gamma = 2\)
  \item
    C. \(\cos^2\alpha + \cos^2\beta + \cos^2\gamma = 3\)
  \item
    D. \(\cos^2\alpha + \cos^2\beta + \cos^2\gamma = 1\)
  \end{itemize}
\item
  设 P, Q 分别为 \(x^{2} + (y - 6)^{2} = 2\) 和椭圆 \(\dfrac{x^{2}}{10} + y^{2} = 1\) 上的点，则 P, Q 两点间的最大距离是

  \begin{itemize}
  \tightlist
  \item
    A. \(5\sqrt{2}\)
  \item
    B. \(\sqrt{46} + \sqrt{2}\)
  \item
    C. \(7 + \sqrt{2}\)
  \item
    D. \(\triangle\sqrt{2}\)
  \end{itemize}
\item
  若存在 \(a \in R\) 且 \(a \neq 0\) , 对任意的 \(x \in R\) , 均有 \(f(x + a) < f(x) + f(a)\) 恒成立, 则称函数 \(f(x)\) 具有性质 \(P\) . 已知: \(q_{1}: f(x)\) 单调递减, 且 \(f(x) > 0\) 恒成立; \(q_{2}: f(x)\) 单调递增, 存在 \(x_{0} < 0\) 使得 \(f(x_{0}) = 0\) , 则是 \(f(x)\) 具有性质 \(P\) 的充分条件是

  \begin{itemize}
  \tightlist
  \item
    A. \(q_{1}\) 和 \(q_{2}\)
  \item
    B. 只有 \(q_{1}\)
  \item
    C. 只有 \(q_{2}\)
  \item
    D. \(q_{1}\) 和 \(q_{2}\) 都不是
  \end{itemize}
\end{enumerate}

\subsection{二、选择题: 大题共 3 小题, 每小题 6 分, 共 18 分. 在每小题给出的四个选项中, 有多项是符合题目要求的, 全部选对的得 6 分, 部分选对的得部分分, 有选错的得 0 分.}\label{ux4e8cux9009ux62e9ux9898-ux5927ux9898ux5171-3-ux5c0fux9898-ux6bcfux5c0fux9898-6-ux5206-ux5171-18-ux5206.-ux5728ux6bcfux5c0fux9898ux7ed9ux51faux7684ux56dbux4e2aux9009ux9879ux4e2d-ux6709ux591aux9879ux662fux7b26ux5408ux9898ux76eeux8981ux6c42ux7684-ux5168ux90e8ux9009ux5bf9ux7684ux5f97-6-ux5206-ux90e8ux5206ux9009ux5bf9ux7684ux5f97ux90e8ux5206ux5206-ux6709ux9009ux9519ux7684ux5f97-0-ux5206.}

\begin{enumerate}
\def\labelenumi{\arabic{enumi}.}
\setcounter{enumi}{8}
\item
  已知曲线 \(C: mx^2 + ny^2 = 1\) ，下列说法正确的是

  \begin{itemize}
  \tightlist
  \item
    A. 若 \(m > n > 0\) ，则 \(C\) 是椭圆，其焦点在 \(y\) 轴上
  \item
    B. 若 \(m = n > 0\) ，则 \(C\) 是圆，其半径为 \(\sqrt{n}\)
  \item
    C. 若 \(mn < 0\) ，则 \(C\) 是双曲线，其渐近线方程为 \(y = \pm \sqrt{-\dfrac{m}{n}} x\)
  \item
    D. 若 \(m = 0, n > 0\) ，则 \(C\) 是两条直线
  \end{itemize}
\end{enumerate}

\ExamQuestionPageBreak

\begin{enumerate}
\def\labelenumi{\arabic{enumi}.}
\setcounter{enumi}{9}
\item
  一组样本数据 \((x_{i},y_{i}), i\in\{1,2,3\cdots,100\}\) . 其中 \(\sum_{i=1}^{100}x_{i}=2\times10^{5},\sum_{i=1}^{100}y_{i}=970\) , 求得其经验回归方程为: \(\hat{y}=-0.02x+a_{1}^{\wedge}\) , 残差为 \(\hat{e}_{i}\) . 对样本数据进行处理: \(x_{i}^{'}=\ln x_{i}\) , 得到新的数据 \((x_{i}^{'},y_{i})\) , 求得其经验回归方程为: \(\hat{y}=-0.42x+a_{2}^{\wedge}\) , 其残差为 \(\hat{u}_{i}^{\wedge}\cdot\hat{e}_{i}^{\wedge},\hat{u}_{i}^{\wedge}\) 分布如图所示, 且 \(\hat{e}\sim N(0,\sigma_{1}^{2}),\hat{u}\sim N(0,\sigma_{2}^{2})\) 则
\end{enumerate}

图1

图2

\pandocbounded{\includegraphics[keepaspectratio]{C:/Mathematics-Knowledge/Secondary-School-Mathematics-Knowledge-Map/临时/测试/images/bcff41cabc8fa738fe173ab60329a2ec398780960e7c1c31f85bb6b00e41a9d2.jpg}}
\pandocbounded{\includegraphics[keepaspectratio]{C:/Mathematics-Knowledge/Secondary-School-Mathematics-Knowledge-Map/临时/测试/images/deba050e6a83a3892e9afd0f848718dfa08f2adf2b35ac2680b17e5ac17c0749.jpg}}

\begin{itemize}
\tightlist
\item
  A. 梯本 \((x_{i},y_{i})\) 负相关
\item
  B. \(a_1^{\wedge} = 49.7\)
\item
  C. \(\sigma_1^2 < \sigma_2^2\)
\item
  D. 处理后的决定系数变大
\end{itemize}

\begin{enumerate}
\def\labelenumi{\arabic{enumi}.}
\setcounter{enumi}{10}
\tightlist
\item
  已知函数 \(f(x) = \ln x - \cos x\) ，将 \(f(x)\) 的所有极值点按照由小到大的顺序排列，得到数列 \(\{x_{n}\}\) ，对于正整数 \(n\) ，则下列说法中正确的有
\end{enumerate}

\begin{itemize}
\tightlist
\item
  A. \((n - \dfrac{1}{2})\pi < x_n < (n + \dfrac{1}{2})\pi\)
\item
  B. \(x_{n+1} - x_n < \pi\)
\item
  C. \(\{ |x_n - n\pi | \}\) 为递减数列
\item
  D. \(f(x_{2n-1}) > \ln \dfrac{(4n-1)\pi}{2}\)
\end{itemize}

\subsection{三、填空题：本题共3小题，每小题5分，共15分.}\label{ux4e09ux586bux7a7aux9898ux672cux9898ux51713ux5c0fux9898ux6bcfux5c0fux98985ux5206ux517115ux5206.}

\begin{enumerate}
\def\labelenumi{\arabic{enumi}.}
\setcounter{enumi}{11}
\item
  记 \(S_{n}\) 为等差数列 \(\{a_{n}\}\) 的前 \(n\) 项和, 若 \(a_{1} = -3, a_{5} = 5\) , 则 \(S_{6} =\) \_\_\_\_.
\item
  袋中有 2 个红球、5 个白球, 现随机不放回地逐个摸出. 设 \(X\) 为``首次出现某一种颜色的球被全部摸出''时所需的摸球次数, 则 \(X = 3\) 的概率为 \_\_\_\_.
\end{enumerate}

\pandocbounded{\includegraphics[keepaspectratio]{C:/Mathematics-Knowledge/Secondary-School-Mathematics-Knowledge-Map/临时/测试/images/1a27997d8ff938c2fb5ae7fb1e48354dc3ab0601ca101650d3cf8f6f930a9724.jpg}}

\begin{enumerate}
\def\labelenumi{\arabic{enumi}.}
\setcounter{enumi}{13}
\item
  一个底面边长和侧棱长均为 4 的正三棱柱密闭容器 ABC - \(A_{1}B_{1}C_{1}\) ，其中盛有一定体积的水，当底面 ABC 水平放置时，水面高为 \(\dfrac{15}{4}\) 。当侧面 \(AA_{1}B_{1}B\) 水平放置时（如图），容器内的水形成新的几何体。若该几何体的所有顶点均在同一个球面上，则该球的表面积为 \_\_\_\_。
\end{enumerate}

\subsection{四、解答题：本题共5小题，共77分．解答应写出文字说明、证明过程或演算步骤.}\label{ux56dbux89e3ux7b54ux9898ux672cux9898ux51715ux5c0fux9898ux517177ux5206ux89e3ux7b54ux5e94ux5199ux51faux6587ux5b57ux8bf4ux660eux8bc1ux660eux8fc7ux7a0bux6216ux6f14ux7b97ux6b65ux9aa4.}

\begin{enumerate}
\def\labelenumi{\arabic{enumi}.}
\setcounter{enumi}{14}
\tightlist
\item
  (13分)
\end{enumerate}

某沙漠地区经过治理,生态系统得到很大改善,野生动物数量有所增加.为调查该地区某种野生动物的数量,将其分成面积相近的200个地块,从这些地块中用简单随机抽样的方法抽取40个作为样区,工作人员将区域内的地块按植被情况分为高覆盖地块、低覆盖地块两类,统计每块样区野生动物数量为``较多''或``较少'',统计结果如下

\begin{longtable}[]{@{}lrrr@{}}
\toprule\noalign{}
& 高覆盖地块 & 低覆盖地块 & 合计 \\
\midrule\noalign{}
\endhead
\bottomrule\noalign{}
\endlastfoot
野生动物数量较多 & 14 & 6 & 20 \\
野生动物数量较少 & 8 & 12 & 20 \\
合计 & 22 & 18 & 40 \\
\end{longtable}


\par\smallskip\noindent（1）依据小概率值 \(\alpha=0.05\) 的独立性检验,判断是否有充分证据认为``该地区此种野生动物数量多少与地块植物覆盖面积高低存在关联'';


\par\smallskip\noindent（2）根据现有统计资料,各地块间植物覆盖面积差异很大.为提高样本的代表性以获得该地区这种野生动物数量更准确的估计,请给出一种你认为更合理的抽样方法,并说明理由.

(注: \(\chi^{2}=\dfrac{n(ad-bc)^{2}}{(a+b)(c+d)(a+c)(b+d)}\) ,其中 \(n=a+b+c+d\) )

\begin{longtable}[]{@{}
  >{\raggedright\arraybackslash}p{(\linewidth - 10\tabcolsep) * \real{0.1304}}
  >{\raggedleft\arraybackslash}p{(\linewidth - 10\tabcolsep) * \real{0.1739}}
  >{\raggedleft\arraybackslash}p{(\linewidth - 10\tabcolsep) * \real{0.1739}}
  >{\raggedleft\arraybackslash}p{(\linewidth - 10\tabcolsep) * \real{0.1739}}
  >{\raggedleft\arraybackslash}p{(\linewidth - 10\tabcolsep) * \real{0.1739}}
  >{\raggedleft\arraybackslash}p{(\linewidth - 10\tabcolsep) * \real{0.1739}}@{}}
\toprule\noalign{}
\begin{minipage}[b]{\linewidth}\raggedright
\(P(\chi^{2} \geqslant k_{\alpha})\)
\end{minipage} & \begin{minipage}[b]{\linewidth}\raggedleft
0.100
\end{minipage} & \begin{minipage}[b]{\linewidth}\raggedleft
0.050
\end{minipage} & \begin{minipage}[b]{\linewidth}\raggedleft
0.010
\end{minipage} & \begin{minipage}[b]{\linewidth}\raggedleft
0.005
\end{minipage} & \begin{minipage}[b]{\linewidth}\raggedleft
0.001
\end{minipage} \\
\midrule\noalign{}
\endhead
\bottomrule\noalign{}
\endlastfoot
\(k_{\alpha}\) & 2.706 & 3.841 & 6.635 & 7.879 & 10.828 \\
\end{longtable}

\ExamQuestionPageBreak

\begin{enumerate}
\def\labelenumi{\arabic{enumi}.}
\setcounter{enumi}{15}
\tightlist
\item
  (15分)
\end{enumerate}

已知函数 \(f(x) = \dfrac{a \cdot 3^x + 2}{3^x + 1} + x^3, a \in R.\)

\begin{enumerate}
\def\labelenumi{(\arabic{enumi})}
\item
  当 \(a = -2\) 时, 求证: 函数 \(f(x)\) 是奇函数;
\item
  当 \(a > -2\) 时, 求 \(f(1) + \dfrac{1}{f(-1)}\) 的最小值.
\end{enumerate}

\begin{enumerate}
\def\labelenumi{\arabic{enumi}.}
\setcounter{enumi}{16}
\tightlist
\item
  (15分)
\end{enumerate}

在 \(\triangle ABC\) 中，角 \(A, B, C\) 的对边分别为 \(a, b, c\) ，已知 \(c = b\tan C\) .

\begin{enumerate}
\def\labelenumi{(\arabic{enumi})}
\item
  若 B 为锐角, 试判断 \(\triangle ABC\) 的形状;
\item
  若 B 为钝角, 求 \(\sin A + \sin C\) 的最大值.
\end{enumerate}

\begin{enumerate}
\def\labelenumi{\arabic{enumi}.}
\setcounter{enumi}{17}
\tightlist
\item
  (17分)
\end{enumerate}

已知函数 \(f(x)=a^{2}x^{2}-3ax\ln x,a>0.\)

\begin{enumerate}
\def\labelenumi{(\arabic{enumi})}
\item
  当 \(a = 1\) 时, 求曲线 \(y = f(x)\) 在 \((1, f(1))\) 处的切线方程;
\item
  讨论 \(f(x)\) 的零点个数；
\item
  当 \(a > 2\) 时, 证明: \(f(x) > 2x\) .
\end{enumerate}

\begin{enumerate}
\def\labelenumi{\arabic{enumi}.}
\setcounter{enumi}{18}
\tightlist
\item
  (17分)
\end{enumerate}

某款游戏预推出一项皮肤抽卡活动,玩家每次抽卡需要花费10元,现有以下两种方案.方案一:没有保底机制,每次抽卡抽中新皮肤的概率为 \(p_{1}\) ;方案二:每次抽卡抽中新皮肤的概率为 \(p_{2}\) ,若连续9次未抽中,则第10次必中新皮肤.已知 \(0<p_{2}<p_{1}<1\) ,玩家按照一、二两种方案进行抽卡,首次抽中新皮肤时的累计花费为X,Y(元).

\begin{enumerate}
\def\labelenumi{(\arabic{enumi})}
\item
  按照方案一, 小明先抽取三次, 没有抽中皮肤, 则他第四次抽中皮肤的概率是多少?
\item
  设第 k 次首次抽中皮肤, 求 \(P(X=10k)\) , \(P(Y=10k)\) ;
\item
  若 \(P_{1} = 2P_{2} = 0.1\) ，根据花费的均值从游戏策划角度选择收益较高的方案。（参考数据： \(0.95^{9} \approx 0.63, 0.95^{10} \approx 0.60\) 。）
\end{enumerate}

\clearpage

\beginExamAnswers
\ExamAnswerSubjectTitle{参考答案}

\subsection{一、选择题}\label{ux4e00ux9009ux62e9ux9898}

\begin{enumerate}
\def\labelenumi{\arabic{enumi}.}
\tightlist
\item
  C．\(-1+\sqrt{3}\mathrm{i}\)\quad{}2. B．\(\{-1,0,1\}\)\quad{}3. B．\(\sqrt2\)\quad{}4. C．\(2\)\quad{}5. D．\(36\)\quad{}6. A．\(\sin^2\alpha+\sin^2\beta+\sin^2\gamma=1\)\quad{}7. \(6\sqrt2\)（按选项位置应为 D，但 D 的 OCR 文本``\(\triangle\sqrt2\)''存在歧义）\quad{}8. A
\end{enumerate}

\subsection{二、选择题}\label{ux4e8cux9009ux62e9ux9898}

\begin{enumerate}
\def\labelenumi{\arabic{enumi}.}
\setcounter{enumi}{8}
\tightlist
\item
  ACD\quad{}10. ABD\quad{}11. ACD
\end{enumerate}

\subsection{三、填空题}\label{ux4e09ux586bux7a7aux9898}

\begin{enumerate}
\def\labelenumi{\arabic{enumi}.}
\setcounter{enumi}{11}
\tightlist
\item
  \(12\)\quad{}13. \(\dfrac{2}{21}\)\quad{}14. \(\dfrac{100\pi}{3}\)
\end{enumerate}

\subsection{四、解答题}\label{ux56dbux89e3ux7b54ux9898}

\begin{enumerate}
\def\labelenumi{\arabic{enumi}.}
\setcounter{enumi}{14}
\tightlist
\item
  （13 分）
\end{enumerate}

\textbf{答案：}
\par\smallskip\noindent（1）没有充分证据认为该地区此种野生动物数量多少与地块植物覆盖面积高低存在关联；
\par\smallskip\noindent（2）可采用分层随机抽样。

\textbf{分析：}


\par\smallskip\noindent（1）根据列联表计算独立性检验统计量，并与显著性水平 \(\alpha=0.05\) 对应的临界值 \(3.841\) 比较；
\par\smallskip\noindent（2）植被覆盖情况是造成地块差异的重要因素，宜先按覆盖情况分层，再在各层中随机抽样。

\textbf{详细解析：}


\par\smallskip\noindent（1）由表中数据取
\[
   a=14,\quad b=6,\quad c=8,\quad d=12,\quad n=40,
   \]
则
\[
   \chi^2
   =\dfrac{40(14\times 12-6\times 8)^2}
   {(14+6)(8+12)(14+8)(6+12)}
   =\dfrac{40\times 120^2}{20\times 20\times 22\times 18}
   =\dfrac{40}{11}
   \approx 3.636.
   \]
因为 \(3.636<3.841\)，

所以在显著性水平 \(\alpha=0.05\) 下不能拒绝``野生动物数量多少与地块植物覆盖面积高低相互独立''的原假设。

因此，没有充分证据认为二者存在关联。


\par\smallskip\noindent（2）可先按植物覆盖面积的高低将全部 \(200\) 个地块分层，再根据各层地块数占全部地块数的比例，在每一层内用简单随机抽样抽取相应数量的地块，使总样本量仍为 \(40\)。

这样能保证覆盖情况不同的地块在样本中均有适当代表；同时，同一层内地块的植被情况相对接近，层内差异较小，通常可以减小抽样误差，

从而提高对该地区野生动物数量估计的准确性。

\ExamQuestionPageBreak

\begin{enumerate}
\def\labelenumi{\arabic{enumi}.}
\setcounter{enumi}{15}
\tightlist
\item
  （15 分）
\end{enumerate}

\textbf{答案：}
\par\smallskip\noindent（1）\(f(x)\) 是奇函数；
\par\smallskip\noindent（2）最小值为 \(2\sqrt3\)。

\textbf{分析：}
\par\smallskip\noindent（1）令 \(a=-2\)，直接计算 \(f(-x)\) 并与 \(-f(x)\) 比较；
\par\smallskip\noindent（2）分别求出 \(f(1)\) 与 \(f(-1)\)，再把所求式化为正数 \(t=a+2\) 的形式，利用基本不等式求最小值。

\textbf{详细解析：}


\par\smallskip\noindent（1）当 \(a=-2\) 时，
\[
   f(x)=\dfrac{-2\cdot 3^x+2}{3^x+1}+x^3.
   \]
函数的定义域为 \(\mathbb R\)，关于原点对称。

又
\[
   \begin{aligned}
   f(-x)
   &=\dfrac{-2\cdot 3^{-x}+2}{3^{-x}+1}-x^3\\
   &=\dfrac{-2+2\cdot 3^x}{1+3^x}-x^3\\
   &=-\left(\dfrac{-2\cdot 3^x+2}{3^x+1}+x^3\right)\\
   &=-f(x).
   \end{aligned}
   \]
因此，函数 \(f(x)\) 是奇函数。


\par\smallskip\noindent（2）由题意 \(a>-2\)。计算得
\[
   f(1)=\dfrac{3a+2}{4}+1=\dfrac{3(a+2)}4,
   \]
\[
   f(-1)=\dfrac{\frac a3+2}{\frac13+1}-1
   =\dfrac{a+6}{4}-1
   =\dfrac{a+2}{4}>0.
   \]
于是
\[
   f(1)+\frac1{f(-1)}
   =\dfrac{3(a+2)}4+\frac4{a+2}.
   \]
令 \(t=a+2>0\)，

则由基本不等式
\[
   \dfrac{3t}{4}+\frac4t
   \geqslant 2\sqrt{\dfrac{3t}{4}\cdot\frac4t}
   =2\sqrt3.
   \]
当且仅当
\[
   \dfrac{3t}{4}=\frac4t,
   \]
即 \(t=\frac4{\sqrt3}\)，也就是 \(a=\frac4{\sqrt3}-2>-2\) 时等号成立。

因此所求最小值为
\[
   \boxed{2\sqrt3}.
   \]

\begin{enumerate}
\def\labelenumi{\arabic{enumi}.}
\setcounter{enumi}{16}
\tightlist
\item
  （15 分）
\end{enumerate}

\textbf{答案：}
\par\smallskip\noindent（1）\(\triangle ABC\) 为以 \(A\) 为直角的直角三角形；
\par\smallskip\noindent（2）\(\sin A+\sin C\) 的最大值为 \(\dfrac98\)。

\textbf{分析：}利用正弦定理把边长关系 \(c=b\tan C\) 转化为角的关系。

根据 \(B\) 是锐角或钝角分别确定 \(B\) 与 \(C\) 的关系；钝角情形下再将目标式化为关于 \(\sin C\) 的二次函数。

\textbf{详细解析：}

由 \(c=b\tan C\) 及 \(b,c>0\) 可知 \(\tan C>0\)，

故 \(0<C<\dfrac{\pi}{2}\)。

由正弦定理，
\[
   \frac cb=\dfrac{\sin C}{\sin B}.
   \]
结合 \(\frac cb=\tan C=\dfrac{\sin C}{\cos C}\)，并注意到 \(\sin C>0\)，

可得
\[
   \sin B=\cos C=\sin\left(\dfrac{\pi}{2}-C\right).
   \]
因此
\[
   B=\dfrac{\pi}{2}-C
   \quad\text{或}\quad
   B=\dfrac{\pi}{2}+C.
   \]


\par\smallskip\noindent（1）若 \(B\) 为锐角，

则只能有
\[
   B=\dfrac{\pi}{2}-C.
   \]
所以 \(B+C=\dfrac{\pi}{2}\)，

从而
\[
   A=\pi-(B+C)=\dfrac{\pi}{2}.
   \]
故 \(\triangle ABC\) 为以 \(A\) 为直角的直角三角形。


\par\smallskip\noindent（2）若 \(B\) 为钝角，

则
\[
   B=\dfrac{\pi}{2}+C,
   \]
从而
\[
   A=\pi-B-C=\dfrac{\pi}{2}-2C.
   \]
由于 \(A>0\)，

故 \(0<C<\dfrac{\pi}{4}\)。于是
\[
   \sin A+\sin C
   =\sin\left(\dfrac{\pi}{2}-2C\right)+\sin C
   =\cos 2C+\sin C.
   \]
令 \(u=\sin C\)，

则 \(0<u<\dfrac{\sqrt2}{2}\)，且
\[
   \cos 2C+\sin C
   =1-2\sin^2C+\sin C
   =1-2u^2+u
   =\frac98-2\left(u-\frac14\right)^2
   \leqslant \frac98.
   \]
当 \(\sin C=\frac14\) 时，确有 \(0<C<\dfrac{\pi}{4}\)，等号可以取得。

因此
\[
   \boxed{\max(\sin A+\sin C)=\frac98}.
   \]

\begin{enumerate}
\def\labelenumi{\arabic{enumi}.}
\setcounter{enumi}{17}
\tightlist
\item
  （17 分）
\end{enumerate}

\textbf{答案：}
\par\smallskip\noindent（1）切线方程为 \(y=-x+2\)；


\par\smallskip\noindent（2）当 \(0<a<\dfrac3e\) 时有两个零点，

当 \(a=\dfrac3e\) 时有一个零点，

当 \(a>\dfrac3e\) 时没有零点；


\par\smallskip\noindent（3）证明见解析。

\textbf{分析：}


\par\smallskip\noindent（1）求导后代入 \(x=1\) 得切线斜率；
\par\smallskip\noindent（2）在定义域 \(x>0\) 内，将零点问题转化为研究 \(g(x)=ax-3\ln x\) 的最小值；
\par\smallskip\noindent（3）把 \(f(x)-2x\) 写成 \(x\) 与另一个函数的乘积，证明后一个因子在其最小点处仍为正。

\textbf{详细解析：}

函数的定义域为 \((0,+\infty)\)。


\par\smallskip\noindent（1）当 \(a=1\) 时，
\[
   f(x)=x^2-3x\ln x,
   \]
所以
\[
   f'(x)=2x-3(\ln x+1).
   \]
又 \(f(1)=1\)，\(f'(1)=2-3=-1\)，

故曲线在 \((1,f(1))=(1,1)\) 处的切线方程为
\[
   y-1=-(x-1),
   \]
即
\[
   \boxed{y=-x+2}.
   \]


\par\smallskip\noindent（2）因为 \(a>0,\ x>0\)，

所以
\[
   f(x)=ax(ax-3\ln x).
   \]
故 \(f(x)=0\) 等价于
\[
   g(x)=ax-3\ln x=0.
   \]
求导得
\[
   g'(x)=a-\frac3x.
   \]
当 \(0<x<\frac3a\) 时，\(g'(x)<0\)；

当 \(x>\frac3a\) 时，\(g'(x)>0\)。

因此 \(g(x)\) 在 \(x=\frac3a\) 处取得最小值
\[
   g\left(\frac3a\right)
   =3-3\ln\frac3a
   =3\ln\dfrac{ae}{3}.
   \]
同时，
\[
   \lim_{x\to0^+}g(x)=+\infty,\qquad
   \lim_{x\to+\infty}g(x)=+\infty.
   \]
因此：

\begin{itemize}
\tightlist
\item
  当 \(0<a<\frac3e\) 时，最小值小于 \(0\)，方程 \(g(x)=0\) 有两个正根，
\end{itemize}

故 \(f(x)\) 有两个零点；
- 当 \(a=\frac3e\) 时，最小值等于 \(0\)，方程有唯一正根 \(x=e\)，

故 \(f(x)\) 有一个零点；
- 当 \(a>\frac3e\) 时，最小值大于 \(0\)，方程无正根，

故 \(f(x)\) 没有零点。


\par\smallskip\noindent（3）当 \(a>2\) 时，
\[
   f(x)-2x
   =x\bigl(a^2x-3a\ln x-2\bigr).
   \]
令
\[
   h(x)=a^2x-3a\ln x-2,\qquad x>0.
   \]
则
\[
   h'(x)=a^2-\dfrac{3a}{x}.
   \]
因此 \(h(x)\) 在 \(x=\frac3a\) 处取得最小值，且
\[
   h\left(\frac3a\right)
   =3a-3a\ln\frac3a-2
   =3a\left(1-\ln\frac3a\right)-2.
   \]
由于 \(a>2\)，

所以 \(0<\frac3a<\frac32\)。利用 \(\ln t<t-1\ (t>0,\ t\ne1)\)（在 \(t=1\) 时也有 \(\ln t\le t-1\)），

可得
\[
   \ln\frac3a\leqslant\frac3a-1<\frac12.
   \]
从而
\[
   h\left(\frac3a\right)
   >3a\cdot\frac12-2
   >1>0.
   \]
故对任意 \(x>0\)，均有 \(h(x)>0\)。

又 \(x>0\)，

所以
\[
   f(x)-2x=xh(x)>0,
   \]
即
\[
   \boxed{f(x)>2x}.
   \]

\begin{enumerate}
\def\labelenumi{\arabic{enumi}.}
\setcounter{enumi}{18}
\tightlist
\item
  （17 分）
\end{enumerate}

\textbf{答案：}
\par\smallskip\noindent（1）\(p_1\)；
\par\smallskip\noindent（2）\(P(X=10k)=(1-p_1)^{k-1}p_1\ (k=1,2,\ldots)\)；对方案二，
\[
   P(Y=10k)=
   \begin{cases}
   (1-p_2)^{k-1}p_2,&1\le k\le9,\\
   (1-p_2)^9,&k=10,\\
   0,&k\ge11;
   \end{cases}
   \]

\par\smallskip\noindent（3）从游戏策划收益角度应选择方案一。

\textbf{分析：}方案一中各次抽卡相互独立，首次抽中次数服从几何分布；方案二中前 \(9\) 次按成功概率 \(p_2\) 独立抽取，第 \(10\) 次在前 \(9\) 次均未抽中时必定抽中。比较两种方案首次抽中时抽卡次数的数学期望即可比较平均累计花费。

\textbf{详细解析：}


\par\smallskip\noindent（1）方案一没有保底机制，每次抽中新皮肤的概率始终为 \(p_1\)，且不受前三次结果影响。

因此，在已知前三次均未抽中的条件下，第四次抽中的概率仍为
\[
   \boxed{p_1}.
   \]


\par\smallskip\noindent（2）对于方案一，第 \(k\) 次首次抽中意味着前 \(k-1\) 次均未抽中且第 \(k\) 次抽中，

所以对 \(k=1,2,\ldots\)，
\[
   \boxed{P(X=10k)=(1-p_1)^{k-1}p_1}.
   \]

对于方案二，

当 \(1\le k\le9\) 时，同理有
\[
   P(Y=10k)=(1-p_2)^{k-1}p_2.
   \]
当 \(k=10\) 时，前 \(9\) 次均未抽中后第 \(10\) 次必定抽中，

故
\[
   P(Y=100)=(1-p_2)^9.
   \]
由于首次抽中最迟发生在第 \(10\) 次，

所以当 \(k\ge11\) 时 \(P(Y=10k)=0\)。综上，
\[
   \boxed{
   P(Y=10k)=
   \begin{cases}
   (1-p_2)^{k-1}p_2,&1\le k\le9,\\
   (1-p_2)^9,&k=10,\\
   0,&k\ge11.
   \end{cases}}
   \]


\par\smallskip\noindent（3）题目末尾写作``\(P_1=2P_2=0.1\)''，而前文定义的是小写 \(p_1,p_2\)。这里按同一概率参数理解，

即
\[
   p_1=0.1,\qquad p_2=0.05.
   \]
设方案一、二中首次抽中的抽卡次数分别为 \(N_1,N_2\)，

则
\[
   E(N_1)=\frac1{p_1}=10,
   \qquad
   E(X)=10E(N_1)=100\text{（元）}.
   \]
对于方案二，令 \(q=1-p_2=0.95\)。

由尾和公式，
\[
   E(N_2)
   =\sum_{k=1}^{10}P(N_2\ge k)
   =\sum_{j=0}^{9}q^j
   =\dfrac{1-q^{10}}{1-q}
   =\dfrac{1-0.95^{10}}{0.05}.
   \]
利用参考数据 \(0.95^{10}\approx0.60\)，

得
\[
   E(N_2)\approx\dfrac{1-0.60}{0.05}=8,
   \qquad
   E(Y)=10E(N_2)\approx80\text{（元）}.
   \]
因为 \(E(X)=100>E(Y)\approx80\)，玩家采用方案一时首次抽中新皮肤前的平均累计花费更高。

因此，从游戏策划收益角度应选择
\[
   \boxed{\text{方案一}}.
   \]

\textbf{注意：}第
\par\smallskip\noindent（3）问的 \(P_1,P_2\) 与题干的 \(p_1,p_2\) 大小写不一致，疑为 OCR 或排版造成的记号差异；上述计算明确按 \(p_1=0.1,\ p_2=0.05\) 解释，未修改源题。


\end{document}
